\section{Conclusão: O Imperativo de uma Abordagem Sistêmica}

Ao longo deste capítulo, mapeou-se um espectro multifacetado de desafios que interagem sinergicamente em múltiplos níveis. As lacunas técnicas corrompem a confiabilidade geométrica e biomecânica do gêmeo digital; as barreiras clínicas geram atrito e desencorajam a adoção pelos profissionais; a letargia regulatória, aliada à ausência de padrões, cria insegurança jurídica e impede a interoperabilidade semântica; as deficiências econômicas esterilizam modelos de negócio viáveis em saúde pública; e as lacunas metodológicas de pesquisa inviabilizam a geração das evidências necessárias para superar todos os entraves anteriores.

A análise aprofundada sugere que é um erro estratégico, e talvez até epistemológico, tratar esses desafios como meros problemas computacionais que serão espontaneamente sanados com o advento de processadores mais potentes ou algoritmos de IA marginalmente melhores. Pelo contrário, as fraturas identificadas são sistêmicas e interdependentes. Protocolos de padronização universal só surgirão com a pressão de marcos regulatórios modernos; as aprovações da ANVISA e CONITEC dependerão intrinsecamente de ensaios clínicos robustos e estudos de custo-efetividade longitudinais, até então inexistentes no âmbito odontológico.

\destaque{A superação das lacunas dos gêmeos digitais não virá de um \textit{breakthrough} tecnológico isolado, mas sim da orquestração de um ecossistema maduro que articule pesquisa clínica, regulamentação ética, interoperabilidade open-source e modelos sustentáveis de financiamento.}

Rompê-lo exigirá uma articulação orquestrada entre o tripé acadêmico, o setor produtivo e o Estado — envolvendo agências de fomento, a ABNT, sociedades científicas odontológicas e os gestores do SUS. Transformar os rigorosos obstáculos aqui expostos em um roteiro pragmático para ação (\textit{roadmap}) não é apenas uma necessidade técnica, mas um imperativo ético para garantir que os gêmeos digitais odontológicos transitem do cativeiro das provas de conceito acadêmicas para a resolução concreta dos agravos bucais da população brasileira.
